% Options for packages loaded elsewhere
\PassOptionsToPackage{unicode}{hyperref}
\PassOptionsToPackage{hyphens}{url}
\PassOptionsToPackage{dvipsnames,svgnames,x11names}{xcolor}
\documentclass[
  11pt,
  letterpaper,
  titlepage]{article}
\usepackage{xcolor}
\usepackage[margin=0.85in]{geometry}
\usepackage{amsmath,amssymb}
\setcounter{secnumdepth}{5}
\usepackage{iftex}
\ifPDFTeX
  \usepackage[T1]{fontenc}
  \usepackage[utf8]{inputenc}
  \usepackage{textcomp} % provide euro and other symbols
\else % if luatex or xetex
  \usepackage{unicode-math} % this also loads fontspec
  \defaultfontfeatures{Scale=MatchLowercase}
  \defaultfontfeatures[\rmfamily]{Ligatures=TeX,Scale=1}
\fi
\usepackage{lmodern}
\ifPDFTeX\else
  % xetex/luatex font selection
\fi
% Use upquote if available, for straight quotes in verbatim environments
\IfFileExists{upquote.sty}{\usepackage{upquote}}{}
\IfFileExists{microtype.sty}{% use microtype if available
  \usepackage[]{microtype}
  \UseMicrotypeSet[protrusion]{basicmath} % disable protrusion for tt fonts
}{}
\usepackage{setspace}
\makeatletter
\@ifundefined{KOMAClassName}{% if non-KOMA class
  \IfFileExists{parskip.sty}{%
    \usepackage{parskip}
  }{% else
    \setlength{\parindent}{0pt}
    \setlength{\parskip}{6pt plus 2pt minus 1pt}}
}{% if KOMA class
  \KOMAoptions{parskip=half}}
\makeatother
\usepackage{longtable,booktabs,array}
\usepackage{caption}
\captionsetup[table]{skip=6pt}
\newcounter{none} % for unnumbered tables
\usepackage{calc} % for calculating minipage widths
% Correct order of tables after \paragraph or \subparagraph
\usepackage{etoolbox}
\makeatletter
\patchcmd\longtable{\par}{\if@noskipsec\mbox{}\fi\par}{}{}
\makeatother
% Allow footnotes in longtable head/foot
\IfFileExists{footnotehyper.sty}{\usepackage{footnotehyper}}{\usepackage{footnote}}
\makesavenoteenv{longtable}
\setlength{\emergencystretch}{3em} % prevent overfull lines
\providecommand{\tightlist}{%
  \setlength{\itemsep}{0pt}\setlength{\parskip}{0pt}}
\usepackage{bookmark}
\IfFileExists{xurl.sty}{\usepackage{xurl}}{} % add URL line breaks if available
\urlstyle{same}
% fallback for those not using the hyperref driver hyperxmp:
\makeatletter
\@ifundefined{xmpquote}{\newcommand{\xmpquote}[1]{#1}}{}
\makeatother
\hypersetup{
  pdftitle={Betterer Prospecting Settings Reference},
  colorlinks=true,
  linkcolor={blue},
  filecolor={Maroon},
  citecolor={Blue},
  urlcolor={blue},
  pdfcreator={LaTeX via pandoc}}

\title{Betterer Prospecting Settings Reference}
\usepackage{etoolbox}
\makeatletter
\providecommand{\subtitle}[1]{% add subtitle to \maketitle
  \apptocmd{\@title}{\par {\large #1 \par}}{}{}
}
\makeatother
\subtitle{Betterer Prospecting 3.4.5 for Vintage Story 1.22.6}
\author{}
\date{2026-08-14}

\begin{document}
\maketitle

{
\hypersetup{linkcolor=black}
\setcounter{tocdepth}{2}
\tableofcontents
}
\setstretch{1.08}
This document describes selected Betterer Prospecting \texttt{3.4.5}
settings as indexed for Vintage Story \texttt{1.22.6}. It is based on
the decompiled mod code and the bundled Config Lib setting manifest.

The setting key \texttt{ReprospectParellelism} is misspelled in the
mod's code and assets. Keep that spelling when editing configuration
files.

\section{Short Answers}\label{short-answers}

{\def\LTcaptype{none} % do not increment counter
\begin{longtable}[]{@{}
  >{\raggedright\arraybackslash}p{(\linewidth - 4\tabcolsep) * \real{0.3333}}
  >{\raggedright\arraybackslash}p{(\linewidth - 4\tabcolsep) * \real{0.3333}}
  >{\raggedright\arraybackslash}p{(\linewidth - 4\tabcolsep) * \real{0.3333}}@{}}
\toprule\noalign{}
\begin{minipage}[b]{\linewidth}\raggedright
Setting
\end{minipage} & \begin{minipage}[b]{\linewidth}\raggedright
Default
\end{minipage} & \begin{minipage}[b]{\linewidth}\raggedright
What It Controls
\end{minipage} \\
\midrule\noalign{}
\endhead
\bottomrule\noalign{}
\endlastfoot
\texttt{ReprospectParellelism} & \texttt{0} & The worker concurrency
used by the \texttt{/btrpr\ reprospect} server command. \\
\texttt{BoreholeRadius} & \texttt{8} & The horizontal radius of the
vertical Borehole scan cylinder. \\
\texttt{BoreholeScansStone} & \texttt{false} & Whether Borehole mode
reports rock/stone blocks in addition to ore. \\
\texttt{BoreholeDmg} & See note below & The amount of prospector pick
durability damage applied by a valid Borehole sample. \\
\end{longtable}
}

\texttt{BoreholeDmg} has a mismatch between two bundled sources. The
decompiled runtime config class initializes it to \texttt{2}, while the
Config Lib setting manifest advertises a default of \texttt{4} with a
range of \texttt{1} to \texttt{10}. Existing saved config values are
authoritative at runtime.

\section{Borehole Mode Context}\label{borehole-mode-context}

Betterer Prospecting replaces the vanilla prospector pick item class and
registers additional tool modes. Borehole mode is one of those modes
when \texttt{AddBoreHoleMode} is enabled.

When a player breaks a block with the prospector pick, Betterer
Prospecting asks the active mode how much durability damage to apply.
The returned value is passed to \texttt{DamageItem}. Borehole mode
normally returns \texttt{BoreholeDmg}.

Borehole mode only accepts samples from the top face of the clicked
block. For a valid top-face sample, it walks a vertical cylinder
centered on the clicked block, from the clicked Y level down to the
bottom of the world. The cylinder radius comes from
\texttt{BoreholeRadius}; the Config Lib manifest defaults that radius to
\texttt{8}, and \texttt{0} makes the scan a straight vertical line.

\section{\texorpdfstring{Matching \texttt{BoreholeRadius} To
\texttt{ProximitySearchRadius\ =\ 5}}{Matching BoreholeRadius To ProximitySearchRadius = 5}}\label{matching-boreholeradius-to-proximitysearchradius-5}

There is no single code-enforced optimum. If ``optimal'' means the
smallest Borehole radius that preserves the same horizontal reach as
\texttt{ProximitySearchRadius\ =\ 5}, use \texttt{BoreholeRadius\ =\ 5}.

This is not the same as the bundled default. Betterer Prospecting ships
\texttt{ProximitySearchRadius\ =\ 5}, but its Config Lib manifest ships
\texttt{BoreholeRadius\ =\ 8}. That default makes Borehole mode a
broader survey tool than Proximity mode.

The reason is the scan shape. \texttt{ProximitySearchRadius} is used by
\texttt{WalkBlocksSphere}, so a radius of \texttt{5} searches a bounded
3D sphere around the clicked block. \texttt{BoreholeRadius} is used by
\texttt{WalkBlocksCylinder}, so a radius of \texttt{5} searches the same
5-block horizontal reach, but repeats that circular footprint from the
clicked Y level down to the bottom of the world.

Integer block positions scanned per horizontal Borehole layer:

{\def\LTcaptype{none} % do not increment counter
\begin{longtable}[]{@{}lrr@{}}
\toprule\noalign{}
\texttt{BoreholeRadius} & Blocks Per Y Layer & Relative To Radius
\texttt{5} \\
\midrule\noalign{}
\endhead
\bottomrule\noalign{}
\endlastfoot
\texttt{5} & \texttt{81} & \texttt{1.00x} \\
\texttt{8} & \texttt{197} & \texttt{2.43x} \\
\texttt{10} & \texttt{317} & \texttt{3.91x} \\
\end{longtable}
}

Because Borehole mode already scans far more depth than Proximity mode,
raising \texttt{BoreholeRadius} above \texttt{5} increases coverage and
scan cost quickly. Use \texttt{8} only when the intent is for Borehole
samples to cover a wider column than Proximity samples. Use \texttt{0}
for a strict vertical line.

\section{\texorpdfstring{\texttt{BoreholeScansStone}}{BoreholeScansStone}}\label{boreholescansstone}

\texttt{BoreholeScansStone} decides whether Borehole mode includes
rock/stone block types in its results.

With the default \texttt{false} value, Borehole mode does not add stone
results. If \texttt{BoreholeScansOre} is still enabled, the scan reports
only qualifying ore blocks found in the vertical cylinder.

When \texttt{BoreholeScansStone} is \texttt{true}, the scan also calls
the mod's rock-detection helper and adds qualifying rock types to the
result list. Ore is checked first, so a block that qualifies as ore is
handled by the ore path rather than the stone path.

If both \texttt{BoreholeScansOre} and \texttt{BoreholeScansStone} are
disabled, Borehole mode does not scan. It sends the
\texttt{borehole-no-filter} message and returns \texttt{1} durability
damage instead of \texttt{BoreholeDmg}.

\section{\texorpdfstring{\texttt{BoreholeDmg}}{BoreholeDmg}}\label{boreholedmg}

\texttt{BoreholeDmg} is the durability cost of a valid Borehole sample.

The Borehole code returns this value after it finishes the scan and
sends the player either the found-results message or the no-results
message. That means a valid top-face Borehole sample consumes
\texttt{BoreholeDmg} durability even when the scan finds no matching ore
or stone.

Two early rejection paths return \texttt{1} instead:

\begin{itemize}
\tightlist
\item
  both ore and stone scanning are disabled;
\item
  the player tries to take the sample from a face other than the top
  face.
\end{itemize}

The bundled English text labels this setting as ``Borehole Damage'' and
describes it as tool damage for Borehole mode. The Config Lib manifest
constrains it to \texttt{1} through \texttt{10}.

\section{\texorpdfstring{\texttt{ReprospectParellelism}}{ReprospectParellelism}}\label{reprospectparellelism}

\texttt{ReprospectParellelism} affects the admin reprospecting command,
not ordinary prospecting pick use.

Betterer Prospecting registers \texttt{/btrpr\ reprospect} behind the
\texttt{controlserver} privilege. The command recalculates stored
prospector-map readings and is described by the mod as an expensive
operation. During that recalculation, the mod uses
\texttt{Parallel.ForEachAsync} over the stored readings for each player
entry.

The setting is used as \texttt{ParallelOptions.MaxDegreeOfParallelism}:

\begin{itemize}
\tightlist
\item
  values greater than \texttt{0} are used directly as the maximum number
  of concurrent reprospecting workers;
\item
  \texttt{0} falls back to \texttt{Environment.ProcessorCount\ /\ 4}.
\end{itemize}

Higher values can finish reprospecting faster, but they increase CPU
load while the command is running. The bundled English text says the
same thing: higher numbers are faster but more CPU-intensive, and
\texttt{0} uses one quarter of the available cores.

The Config Lib manifest allows \texttt{0} through \texttt{100}. For
normal servers, \texttt{0} is the conservative choice. Increase it only
for planned maintenance, preferably while few players are online.

\section{Notes And Quirks}\label{notes-and-quirks}

\begin{itemize}
\tightlist
\item
  The spelling \texttt{Parellelism} appears in the actual key and UI
  translation. Do not silently correct it to \texttt{Parallelism} in
  config files.
\item
  The decompiled \texttt{/btrpr\ reprospect} command accepts an optional
  player argument, but the visible \texttt{3.4.5} method body loops
  through the stored map-layer readings collection as a whole. Verify
  in-game before relying on the command for per-player-only
  reprospecting.
\item
  The Config Lib manifest and the direct \texttt{ModConfig} class
  disagree on \texttt{BoreholeDmg}'s initial default. Check the saved
  runtime config when diagnosing a server's actual Borehole durability
  cost.
\end{itemize}

\section{Source Paths}\label{source-paths}

All paths below are under
\texttt{versions/1.22.6/mods/bettererprospecting\_3.4.5}.

Under
\texttt{decompiled\ assembly/BetterErProspecting/BetterErProspecting/}:

\begin{itemize}
\tightlist
\item
  \texttt{Config/ModConfig.cs}
\item
  \texttt{Item/ItemBetterErProspectingPick.cs}
\item
  \texttt{Prospecting/ProspectingSystem.cs}
\end{itemize}

Under \texttt{assets/bettererprospecting/}:

\begin{itemize}
\tightlist
\item
  \texttt{config/configlib-patches.json}
\item
  \texttt{lang/en.json}
\end{itemize}

\end{document}
